% ========================================================
% BÀI TẬP TOÁN 5 - KẾT NỐI TRI THỨC VỚI CUỘC SỐNG
% Bài 22: Phép chia số thập phân
% Định dạng: BTTeX Part III - Trắc nghiệm điền số ngắn
% ========================================================

% Câu 1
\begin{ex}
  Tính thương của phép chia: $24 : 1{,}5$.
  \shortans{16}
  \loigiai{
    $24 : 1{,}5 = 240 : 15 = 16$.
  }
\end{ex}

% Câu 2
\begin{ex}
  Tính thương: $45 : 2{,}5$.
  \shortans{18}
  \loigiai{
    $45 : 2{,}5 = 450 : 25 = 18$.
  }
\end{ex}

% Câu 3
\begin{ex}
  Tính thương: $75 : 0{,}5$.
  \shortans{150}
  \loigiai{
    $75 : 0{,}5 = 75 \times 2 = 150$.
  }
\end{ex}

% Câu 4
\begin{ex}
  Tính thương: $18 : 0{,}25$.
  \shortans{72}
  \loigiai{
    $18 : 0{,}25 = 18 \times 4 = 72$.
  }
\end{ex}

% Câu 5
\begin{ex}
  Tính thương: $12{,}5 : 0{,}5$.
  \shortans{25}
  \loigiai{
    $12{,}5 \times 2 = 25$.
  }
\end{ex}

% Câu 6
\begin{ex}
  Tính: $14{,}4 : 1{,}2$.
  \shortans{12}
  \loigiai{
    $144 : 12 = 12$.
  }
\end{ex}

% Câu 7
\begin{ex}
  Tính: $48{,}6 : 9$. Lấy kết quả nhân với 10 được bao nhiêu?
  \shortans{54}
  \loigiai{
    $48{,}6 : 9 = 5{,}4$. Nhân 10 được 54.
  }
\end{ex}

% Câu 8
\begin{ex}
  Tính thương: $100 : 12{,}5$.
  \shortans{8}
  \loigiai{
    $100 : 12{,}5 = 8$.
  }
\end{ex}

% Câu 9
\begin{ex}
  Tìm $x$ biết: $x \times 1{,}5 = 45$.
  \shortans{30}
  \loigiai{
    $x = 45 : 1{,}5 = 30$.
  }
\end{ex}

% Câu 10
\begin{ex}
  Tìm $x$ biết: $x \times 2{,}5 = 100$.
  \shortans{40}
  \loigiai{
    $x = 100 : 2{,}5 = 40$.
  }
\end{ex}

% Câu 11
\begin{ex}
  Một mảnh vườn hình chữ nhật có diện tích $96\text{ m}^2$, chiều rộng là $6{,}4\text{ m}$. Chiều dài mảnh vườn là bao nhiêu mét?
  \shortans{15}
  \loigiai{
    Chiều dài là: $96 : 6{,}4 = 15\text{ m}$.
  }
\end{ex}

% Câu 12
\begin{ex}
  Có $45$ lít dầu hỏa đem chia đều vào các chai, mỗi chai chứa $0{,}75$ lít. Chia được tất cả bao nhiêu chai?
  \shortans{60}
  \loigiai{
    Số chai là: $45 : 0{,}75 = 60$ (chai).
  }
\end{ex}

% Câu 13
\begin{ex}
  Một người may $8$ bộ quần áo hết $18\text{ m}$ vải. Hỏi may $12$ bộ quần áo như thế hết bao nhiêu mét vải?
  \shortans{27}
  \loigiai{
    1 bộ hết $18 : 8 = 2{,}25\text{ m}$. 12 bộ hết $2{,}25 \times 12 = 27\text{ m}$.
  }
\end{ex}

% Câu 14
\begin{ex}
  Một ô tô đi quãng đường $135\text{ km}$ hết $2{,}5\text{ giờ}$. Vận tốc của ô tô là bao nhiêu ki-lô-mét trên giờ?
  \shortans{54}
  \loigiai{
    Vận tốc là: $135 : 2{,}5 = 54\text{ km/h}$.
  }
\end{ex}

% Câu 15
\begin{ex}
  Biết $2{,}5$ lít dầu cân nặng $2\text{ kg}$. Hỏi $10$ lít dầu như thế cân nặng bao nhiêu ki-lô-gam?
  \shortans{8}
  \loigiai{
    1 lít nặng $2 : 2{,}5 = 0{,}8\text{ kg}$. 10 lít nặng $0{,}8 \times 10 = 8\text{ kg}$.
  }
\end{ex}

% Câu 16
\begin{ex}
  Tính giá trị biểu thức: $120 : 0{,}5 : 4$.
  \shortans{60}
  \loigiai{
    $120 : 0{,}5 = 240$. $240 : 4 = 60$.
  }
\end{ex}

% Câu 17
\begin{ex}
  Tính giá trị: $(36{,}5 + 13{,}5) : 2{,}5$.
  \shortans{20}
  \loigiai{
    $50 : 2{,}5 = 20$.
  }
\end{ex}

% Câu 18
\begin{ex}
  Tìm số tự nhiên $y$ biết: $75 : y = 1{,}25$.
  \shortans{60}
  \loigiai{
    $y = 75 : 1{,}25 = 60$.
  }
\end{ex}

% Câu 19
\begin{ex}
  Một thanh sắt dài $1{,}6\text{ m}$ cân nặng $16\text{ kg}$. Hỏi thanh sắt cùng loại dài $2{,}5\text{ m}$ cân nặng bao nhiêu ki-lô-gam?
  \shortans{25}
  \loigiai{
    1 m nặng $16 : 1{,}6 = 10\text{ kg}$. $2{,}5\text{ m}$ nặng $10 \times 2{,}5 = 25\text{ kg}$.
  }
\end{ex}

% Câu 20
\begin{ex}
  Người ta rót $30$ lít mật ong vào các chai loại $0{,}5$ lít. Rót được bao nhiêu chai?
  \shortans{60}
  \loigiai{
    $30 : 0{,}5 = 60$ (chai).
  }
\end{ex}

% Câu 21
\begin{ex}
  Tính thương: $200 : 0{,}8$.
  \shortans{250}
  \loigiai{
    $200 : 0{,}8 = 250$.
  }
\end{ex}

% Câu 22
\begin{ex}
  Một hình tam giác có diện tích $30\text{ cm}^2$, chiều cao là $4\text{ cm}$. Độ dài cạnh đáy tương ứng là bao nhiêu xăng-ti-mét?
  \shortans{15}
  \loigiai{
    Đáy là $(30 \times 2) : 4 = 15\text{ cm}$.
  }
\end{ex}

% Câu 23
\begin{ex}
  Tính giá trị biểu thức: $15 : 1{,}5 + 24 : 1{,}2$.
  \shortans{30}
  \loigiai{
    $10 + 20 = 30$.
  }
\end{ex}

% Câu 24
\begin{ex}
  Một tổ công nhân đắp được $48\text{ m}$ đường trong $3{,}2$ ngày. Hỏi mỗi ngày tổ đó đắp được bao nhiêu mét đường?
  \shortans{15}
  \loigiai{
    $48 : 3{,}2 = 15\text{ m}$.
  }
\end{ex}

% Câu 25
\begin{ex}
  Tìm số tự nhiên $x$ biết $x \times 0{,}4 = 18 : 1{,}5$.
  \shortans{30}
  \loigiai{
    $18 : 1{,}5 = 12$. $x = 12 : 0{,}4 = 30$.
  }
\end{ex}
